Weer goot de maniak
iets in Emmy’s drinkbak
verstopt in haar schortzak.

Dit keer goot de jager,
als haar waterdrager
en wellicht aanklager,

het niet uit in de tuin.
Die plek in de moestuin
kleurde immers snel bruin.

Dit keer deed de wreker
het in Jeanne’s beker.
Zo wist hij trefzeker

of Jeanne echt Emmy,
zonder zijn permissie,
als geheime missie,

wilde vergiftigen.
Haar nijd bevestigen.
Zijn eigen ledigen.

Emmy dronk als hondje
met lebberend mondje
en wiebelend kontje

naakt geknield op de vloer.
Hij draaide haar een loer
door deze in-huis-hoer

stevig van achteren
te overmeesteren.
Jeanne te treiteren

door haar af te leiden.
De indruk te mijden
dat de dood zal scheiden.

Jeanne voelde woede,
was niet op haar hoede
en dat beïnvloedde

sterk haar opmerkzaamheid.
Ze dronk hetgeen bereid.
Zijn opgekropte nijd

spoot hij in de thuisslet.
Aldus niet opgelet,
lag ze in haar sterfbed

een uur of twee later.
Jeanne zag haar flater
als een flinke kater.

Woede vertaald in wraak
stelde haar aan de kaak.
Haar mislukte moordzaak

had zich zo openbaart.
De lucht was geklaard,
en ze kon zo bedaard

op haar woede ingaan.
“Ik ben ook vreemdgegaan.
Ik legde het ooit aan

met onze heilsprofeet,
toen jij al lang en breed
het steeds met Svenska deed.

De profeet onthulde,
hoe jij je dag vulde.
Wat ik sindsdien dulde,

omdat de leer ons zegt;
wie een ander berecht,
wordt zelf eveneens slecht.

Ik herstelde bijkans
de verstoorde balans,
middels mijn eigen sjans

met onze heilsprofeet.
Toch bleef de diepe leed,
die je me ooit aandeed.

Je schond je gelofte!”
De jager ontplofte,
waardoor Jeanne bofte

dat ze toch al dood ging.
Een eigen onthulling
was zijn verdediging.

Zo biechtte de jager,
dat de gezagsdrager,
als ook de aanklager

van Svenska’s rooftochten,
hem bedrukt opzochten
op zijn kroegentochten.

Eerst wilde hij beiden,
zonder medelijden,
hun kelen doorsnijden.

In zijn dronkentoestand
kon hij wel met verstand
luisteren naar hun kant.

“Spijts zijn eigen dogma,
met moord als dillema,
heeft de profeet Svenska

aan de Bisschop verlinkt.
Dat misdaadzaakje stinkt,
want er werd afgevinkt

iemand van zijn volglijst.
Wat volgens ons bewijst
dat schouwing is vereist.

Want er zijn geruchten
dat zij niet kon duchten
en ook wilde luchten

dat hij foute dingen
met zijn volgelingen
slim wist af te dwingen.

Hij scheen te worstelen
met de apostelen.
Hij kon gladborstelen

Svenska’s beschuldiging
door haar terechtstelling
voor haar openbaring.

Wij zijn maar speelmannen
van blijkbaar zijn plannen
en daarvoor verbannen.”

De jager wist genoeg.
Nadat hij hen dood sloeg,
verliet hij zijn stamkroeg

en zocht toen de bisschop
in zijn spreekkamer op.
Bang voor zijn boevenkop

bevestigde die snel
het moordende zwijgspel
en deed schuw een voorstel.

De jager kreeg zijn woord
dat het lek werd vermoord.
Hij vond het ongehoord

dat de sekteleider
als geestrijk bevrijder
deze seksuitglijder

wis op zijn naam had staan.
Hij had hem met argwaan
eerst zijn gang laten gaan.

Met de jagers dreiging
koos hij voor bedekking
door zijn verwijdering.

“Je profeet neukte meer”,
gaf de jager als sneer.
Hij keek op Jeanne neer.

Het leek haar te deren.
Toch wist ze te keren
door zelf te beweren:

“Oen, je bent onvruchtbaar.
Ja, uilskuiken, lach maar.
Maar het is toch echt waar.

Ik heb je bedrogen
en je voorgelogen.
Je zal een kind zogen

dat, net als bij je broer,
is verwekt bij een hoer.
Want je dreutelt wel stoer,

maar zaadjes zijn er niet.
Een andere naaipiet
gaf die aan je rotgriet.”

Deze laatste woorden,
zonder nog antwoorden
om te verantwoorden

deze pijnlijke klucht,
namen haar laatste lucht.
Jeanne ging met een zucht.
